
\documentclass[11pt]{article}

\usepackage[letterpaper,margin=1in]{geometry}
\usepackage[T1]{fontenc}
\usepackage{lmodern}
\usepackage{microtype}
\usepackage{xcolor}
\usepackage{listings}
\usepackage{longtable}
\usepackage{array}
\usepackage{booktabs}
\usepackage{enumitem}
\usepackage{hyperref}

\hypersetup{
    colorlinks=true,
    linkcolor=blue,
    urlcolor=blue
}

\definecolor{codegray}{rgb}{0.96,0.96,0.96}

\lstset{
    backgroundcolor=\color{codegray},
    basicstyle=\ttfamily\small,
    columns=fullflexible,
    keepspaces=true,
    breaklines=true,
    frame=single,
    showstringspaces=false,
    tabsize=2
}

\setlist{nosep}

\title{
    \textbf{SciVault Question Format}\\
    Human-Readable Specification\\
    \large SciVault-QF Version 0.1
}

\author{SciVault}
\date{}

\begin{document}

\maketitle

\begin{abstract}
The SciVault Question Format (SciVault-QF) is an open,
application-independent JSON format for representing educational questions
and question banks. It provides a common question-content layer for the
SciVault software suite and may also be implemented by third-party
applications.

SciVault-QF separates question content from presentation, assessment
sessions, student records, and results. This allows a question to be
authored once and reused for test generation, peer instruction, practice,
homework, online assessment, analytics, and other instructional purposes.
\end{abstract}

\tableofcontents
\newpage

% ============================================================
\section{Scope}
% ============================================================

SciVault-QF defines the representation of educational questions and
collections of questions.

It is intended to support applications including:

\begin{itemize}
    \item question-bank editors;
    \item test and worksheet generators;
    \item peer-instruction systems;
    \item classroom-response systems;
    \item homework and practice systems;
    \item online assessments;
    \item item-analysis systems;
    \item learning analytics;
    \item learning-management-system integrations; and
    \item compatible third-party educational software.
\end{itemize}

SciVault-QF does \textbf{not} define student accounts, classroom rosters,
assessment sessions, grades, or student responses.

% ============================================================
\section{Design Goals}
% ============================================================

SciVault-QF is designed to be:

\begin{itemize}
    \item human-readable;
    \item machine-readable;
    \item portable;
    \item extensible;
    \item suitable for version control;
    \item independent of any particular database;
    \item independent of any particular SciVault application;
    \item suitable for mathematical and scientific content;
    \item capable of referencing media;
    \item suitable for static and dynamically generated questions; and
    \item usable by software outside the SciVault suite.
\end{itemize}

% ============================================================
\section{Format Identity}
% ============================================================

The official name of the format is:

\begin{center}
\textbf{SciVault Question Format}
\end{center}

The abbreviated machine-readable identifier is:

\begin{lstlisting}
SciVault-QF
\end{lstlisting}

Version 0.1 files shall contain:

\begin{lstlisting}
"format": "SciVault-QF",
"format_version": "0.1"
\end{lstlisting}

% ============================================================
\section{File Representation}
% ============================================================

SciVault-QF version 0.1 uses JSON.

\begin{itemize}
    \item \textbf{Encoding:} UTF-8
    \item \textbf{Extension:} \texttt{.json}
    \item \textbf{Recommended MIME type:} \texttt{application/json}
\end{itemize}

Version 0.1 does not define a proprietary file extension.

% ------------------------------------------------------------
\subsection{Future Packages}
% ------------------------------------------------------------

A future specification may define a packaged question-bank format containing
a SciVault-QF document together with images, audio, video, or other
resources.

SciVault-QF 0.1 does not specify the archive technology or package
extension.

ZIP, TAR, TAR.GZ, and other container formats are therefore neither required
nor prohibited by this specification.

% ============================================================
\section{Document Structure}
% ============================================================

A SciVault-QF document contains metadata and one or more questions.

\begin{lstlisting}
{
  "format": "SciVault-QF",
  "format_version": "0.1",

  "metadata": {
    ...
  },

  "questions": [
    ...
  ]
}
\end{lstlisting}

The top-level fields are:

\begin{longtable}{|>{\ttfamily}p{3.0cm}|p{2.0cm}|p{1.7cm}|p{7.5cm}|}
\hline
\textbf{Field} & \textbf{Type} & \textbf{Required} &
\textbf{Description} \\ \hline
\endfirsthead

format & String & Yes &
Must contain \texttt{SciVault-QF}. \\ \hline

format\_version & String & Yes &
Version of this specification. \\ \hline

metadata & Object & Yes &
Metadata describing the question collection. \\ \hline

questions & Array & Yes &
Array containing one or more questions. \\ \hline

extensions & Object & No &
Namespaced extension data. \\ \hline

\end{longtable}

Unknown fields shall not be inserted directly into core SciVault-QF
objects. Additional data shall use the extension mechanism described later
in this specification.

% ============================================================
\section{Question-Set Metadata}
% ============================================================

Example:

\begin{lstlisting}
"metadata": {
  "id": "org.scivault.physics.kinematics",
  "title": "Kinematics Concept Questions",
  "description":
    "Conceptual and quantitative introductory kinematics questions.",
  "author": "Jane Doe",
  "license": "CC-BY-4.0",
  "version": "1.0",
  "language": "en-US",
  "tags": [
    "physics",
    "kinematics"
  ],
  "created": "2026-08-19T10:30:00-06:00",
  "modified": "2026-08-19T10:45:00-06:00"
}
\end{lstlisting}

\begin{longtable}{|>{\ttfamily}p{2.8cm}|p{1.8cm}|p{1.7cm}|p{7.8cm}|}
\hline
\textbf{Field} & \textbf{Type} & \textbf{Required} &
\textbf{Description} \\ \hline
\endfirsthead

id & String & Recommended &
Stable identifier for the question collection. \\ \hline

title & String & Yes &
Human-readable title. \\ \hline

description & String & No &
Description of the collection. \\ \hline

author & String & No &
Author or responsible organization. \\ \hline

license & String & Yes &
License governing the question content. \\ \hline

version & String & No &
Version of the question collection. \\ \hline

language & String & No &
Content language or locale, preferably expressed using a BCP~47 tag such
as \texttt{en-US}. \\ \hline

tags & Array & No &
Descriptive keywords applying to the collection. \\ \hline

created & String & No &
Creation timestamp in ISO~8601 date-time format. \\ \hline

modified & String & No &
Modification timestamp in ISO~8601 date-time format. \\ \hline

extensions & Object & No &
Namespaced extension data. \\ \hline

\end{longtable}

The number of questions shall not be stored separately. Applications shall
determine it from the length of the \texttt{questions} array.

% ============================================================
\section{Identifiers}
% ============================================================

Each question shall have an ID unique within its question collection.

Example:

\begin{lstlisting}
"id": "q001"
\end{lstlisting}

A stable namespace-style identifier is recommended for the question set:

\begin{lstlisting}
"id": "org.scivault.physics.kinematics"
\end{lstlisting}

The combination:

\begin{lstlisting}
org.scivault.physics.kinematics:q001
\end{lstlisting}

may therefore be treated as a globally meaningful reference to the
question.

Applications shall not assume that IDs are numeric or sequential.

Answer choices shall also have stable IDs unique within their question.

% ============================================================
\section{Common Question Structure}
% ============================================================

All questions shall contain:

\begin{itemize}
    \item \texttt{id};
    \item \texttt{type};
    \item \texttt{question}; and
    \item \texttt{answer}.
\end{itemize}

Questions may additionally contain:

\begin{itemize}
    \item \texttt{choices};
    \item \texttt{explanation};
    \item \texttt{hint};
    \item \texttt{difficulty};
    \item \texttt{tags};
    \item \texttt{objectives};
    \item \texttt{media};
    \item \texttt{variables};
    \item \texttt{settings};
    \item \texttt{metadata}; and
    \item \texttt{extensions}.
\end{itemize}

Example:

\begin{lstlisting}
{
  "id": "q001",
  "type": "multiple_choice",
  "question": "What is $2+2$?",
  "choices": [
    {"id": "A", "text": "3"},
    {"id": "B", "text": "4"},
    {"id": "C", "text": "5"}
  ],
  "answer": "B",
  "explanation": "$2+2=4$.",
  "difficulty": 1,
  "tags": ["math", "arithmetic"]
}
\end{lstlisting}

% ============================================================
\section{Text and Mathematical Notation}
% ============================================================

Question prompts, choices, hints, explanations, and other textual content
may contain LaTeX-compatible mathematical notation.

Example:

\begin{lstlisting}
"What is the kinetic energy of a $4\\,\\mathrm{kg}$ object
moving at $3\\,\\mathrm{m/s}$?"
\end{lstlisting}

Because backslashes have special meaning in JSON, LaTeX backslashes shall
be escaped.

For example:

\begin{lstlisting}
\frac{1}{2}mv^2
\end{lstlisting}

is represented within a JSON string as:

\begin{lstlisting}
"\\frac{1}{2}mv^2"
\end{lstlisting}

Applications should use a mathematical rendering system such as MathJax or
KaTeX.

% ============================================================
\section{Difficulty}
% ============================================================

The optional \texttt{difficulty} field shall contain an integer from
0 through 10 inclusive.

The scale is intended to permit question authors and applications to
classify questions with greater resolution than a small categorical scale.

The endpoints may generally be interpreted as:

\begin{center}
\begin{tabular}{|c|l|}
\hline
0 & Trivial or introductory \\ \hline
5 & Moderate \\ \hline
10 & Extremely difficult \\ \hline
\end{tabular}
\end{center}

Intermediate values represent increasing levels of difficulty.

Difficulty is primarily intended for:

\begin{itemize}
    \item filtering;
    \item question-bank organization;
    \item automatic test generation;
    \item differentiated practice;
    \item adaptive systems; and
    \item instructional planning.
\end{itemize}

A difficulty value is an author or system estimate unless supported by
empirical item statistics.

% ============================================================
\section{Tags and Learning Objectives}
% ============================================================

Tags provide descriptive classification.

Example:

\begin{lstlisting}
"tags": [
  "physics",
  "kinematics",
  "constant-acceleration"
]
\end{lstlisting}

Learning objectives are stored separately:

\begin{lstlisting}
"objectives": [
  "KIN-VELOCITY",
  "PHY.KIN.3"
]
\end{lstlisting}

This distinction permits questions to be searched broadly by topic while
also associating them with specific assessable outcomes.

SciVault-QF does not require a particular learning-objective or standards
framework.

% ============================================================
\section{Multiple-Choice Questions}
% ============================================================

Multiple-choice questions use:

\begin{lstlisting}
"type": "multiple_choice"
\end{lstlisting}

They shall contain two or more choices and exactly one correct answer.

\begin{lstlisting}
{
  "id": "q001",
  "type": "multiple_choice",

  "question":
    "An object starts from rest and accelerates at "
    "$2\\,\\mathrm{m/s^2}$ for $5\\,\\mathrm{s}$. "
    "What is its final velocity?",

  "choices": [
    {
      "id": "A",
      "text": "$5\\,\\mathrm{m/s}$"
    },
    {
      "id": "B",
      "text": "$10\\,\\mathrm{m/s}$"
    },
    {
      "id": "C",
      "text": "$15\\,\\mathrm{m/s}$"
    },
    {
      "id": "D",
      "text": "$20\\,\\mathrm{m/s}$"
    }
  ],

  "answer": "B",

  "explanation":
    "Using $v=v_0+at$, "
    "$v=0+(2)(5)=10\\,\\mathrm{m/s}$.",

  "difficulty": 2
}
\end{lstlisting}

The \texttt{answer} shall reference the stable ID of the correct choice.

Displayed answer text shall not serve as the answer identifier.

% ============================================================
\section{Multiple-Select Questions}
% ============================================================

Questions having more than one correct choice use:

\begin{lstlisting}
"type": "multiple_select"
\end{lstlisting}

The answer shall contain an array of correct choice IDs.

\begin{lstlisting}
{
  "id": "q002",
  "type": "multiple_select",

  "question":
    "Which of the following are vector quantities?",

  "choices": [
    {"id": "A", "text": "Velocity"},
    {"id": "B", "text": "Mass"},
    {"id": "C", "text": "Acceleration"},
    {"id": "D", "text": "Temperature"}
  ],

  "answer": [
    "A",
    "C"
  ]
}
\end{lstlisting}

% ============================================================
\section{True/False Questions}
% ============================================================

True/false questions use:

\begin{lstlisting}
"type": "true_false"
\end{lstlisting}

They shall contain exactly two choices with the IDs:

\begin{lstlisting}
true
false
\end{lstlisting}

Example:

\begin{lstlisting}
{
  "id": "q003",
  "type": "true_false",

  "question":
    "An object moving at constant velocity has zero acceleration.",

  "choices": [
    {"id": "true", "text": "True"},
    {"id": "false", "text": "False"}
  ],

  "answer": "true"
}
\end{lstlisting}

% ============================================================
\section{Short-Answer Questions}
% ============================================================

Short-answer questions use:

\begin{lstlisting}
"type": "short_answer"
\end{lstlisting}

The answer may contain one acceptable response:

\begin{lstlisting}
"answer": "photosynthesis"
\end{lstlisting}

or multiple acceptable textual responses:

\begin{lstlisting}
"answer": [
  "photosynthesis",
  "Photosynthesis"
]
\end{lstlisting}

Applications may implement normalization, case-insensitive comparison,
teacher review, or other appropriate grading behavior.

% ============================================================
\section{Fill-in-the-Blank Questions}
% ============================================================

Fill-in-the-blank questions use:

\begin{lstlisting}
"type": "fill_in_the_blank"
\end{lstlisting}

Version 0.1 reserves:

\begin{lstlisting}
{{blank}}
\end{lstlisting}

as the blank placeholder.

Example:

\begin{lstlisting}
{
  "id": "q004",
  "type": "fill_in_the_blank",

  "question":
    "The SI unit of force is the {{blank}}.",

  "answer": [
    "newton",
    "Newton",
    "N"
  ]
}
\end{lstlisting}

Version 0.1 defines one scored blank per question.

Multiple independently scored blanks are reserved for a future revision.

% ============================================================
\section{Numerical Questions}
% ============================================================

Numerical questions use:

\begin{lstlisting}
"type": "numerical"
\end{lstlisting}

A static numerical answer may contain:

\begin{lstlisting}
"answer": {
  "value": 20,
  "tolerance": 0.1,
  "unit": "m/s"
}
\end{lstlisting}

The \texttt{value} is the expected answer.

The optional \texttt{tolerance} specifies an absolute numerical tolerance.

The optional \texttt{unit} specifies the expected physical unit.

A dynamically calculated answer may instead use an expression:

\begin{lstlisting}
"answer": {
  "expression": "a * t",
  "tolerance": 0.01,
  "unit": "m/s"
}
\end{lstlisting}

Exactly one of \texttt{value} or \texttt{expression} shall be present.

% ============================================================
\section{Hints and Explanations}
% ============================================================

The optional \texttt{hint} field contains assistance that may be displayed
before the answer is revealed.

Example:

\begin{lstlisting}
"hint":
  "Identify the initial velocity, acceleration, and elapsed time."
\end{lstlisting}

The optional \texttt{explanation} field contains feedback, reasoning, or a
worked solution.

Example:

\begin{lstlisting}
"explanation":
  "Using $v=v_0+at$, "
  "$v=0+(2)(5)=10\\,\\mathrm{m/s}$."
\end{lstlisting}

Consuming applications determine when hints and explanations are displayed.

% ============================================================
\section{Media}
% ============================================================

Questions and answer choices may reference external media.

Example:

\begin{lstlisting}
"media": [
  {
    "id": "diagram1",
    "type": "image",
    "src": "media/incline.svg",
    "alt":
      "A block resting on a 30 degree inclined plane."
  }
]
\end{lstlisting}

Version 0.1 recognizes the media types:

\begin{itemize}
    \item \texttt{image};
    \item \texttt{audio};
    \item \texttt{video}.
\end{itemize}

Media objects may contain:

\begin{itemize}
    \item \texttt{id};
    \item \texttt{type};
    \item \texttt{src};
    \item \texttt{alt};
    \item \texttt{caption};
    \item \texttt{credit};
    \item \texttt{license}; and
    \item \texttt{extensions}.
\end{itemize}

Images shall contain an \texttt{alt} field.

Relative paths are recommended for locally distributed resources.

% ============================================================
\section{Dynamic Questions}
% ============================================================

Dynamic questions may define variables whose values are generated when the
question is instantiated.

Variables are referenced using double braces.

Example:

\begin{lstlisting}
"question":
  "An object accelerates at {{a}} m/s^2 "
  "for {{t}} s. What is its change in velocity?"
\end{lstlisting}

Version 0.1 defines:

\begin{itemize}
    \item integer variables;
    \item decimal variables;
    \item choice variables.
\end{itemize}

% ------------------------------------------------------------
\subsection{Integer Variables}
% ------------------------------------------------------------

\begin{lstlisting}
"a": {
  "type": "integer",
  "min": 1,
  "max": 5
}
\end{lstlisting}

An optional positive \texttt{step} may be specified.

% ------------------------------------------------------------
\subsection{Decimal Variables}
% ------------------------------------------------------------

\begin{lstlisting}
"mu": {
  "type": "decimal",
  "min": 0.1,
  "max": 1.0,
  "step": 0.1
}
\end{lstlisting}

% ------------------------------------------------------------
\subsection{Choice Variables}
% ------------------------------------------------------------

\begin{lstlisting}
"planet": {
  "type": "choice",
  "values": [
    "Earth",
    "Mars",
    "Moon"
  ]
}
\end{lstlisting}

Choice-variable values may be strings, numbers, or Boolean values.

% ============================================================
\section{Dynamic Expressions}
% ============================================================

A dynamic numerical answer may be calculated from variables.

Example:

\begin{lstlisting}
{
  "id": "q005",
  "type": "numerical",

  "question":
    "An object starts from rest and accelerates at "
    "{{a}} m/s^2 for {{t}} s. "
    "What is its final speed?",

  "variables": {
    "a": {
      "type": "integer",
      "min": 1,
      "max": 5
    },

    "t": {
      "type": "integer",
      "min": 2,
      "max": 10
    }
  },

  "answer": {
    "expression": "a * t",
    "tolerance": 0.01,
    "unit": "m/s"
  }
}
\end{lstlisting}

Expressions shall be evaluated by a restricted mathematical expression
engine.

SciVault-QF files shall never cause arbitrary programming-language code to
be executed.

The complete expression language is reserved for further specification.

% ============================================================
\section{Randomization}
% ============================================================

Applications may randomize:

\begin{itemize}
    \item question selection;
    \item question order;
    \item dynamic variable values; and
    \item answer-choice order.
\end{itemize}

Correctness is associated with stable choice IDs rather than displayed
positions.

A question may recommend answer-choice shuffling:

\begin{lstlisting}
"settings": {
  "shuffle_choices": true
}
\end{lstlisting}

% ============================================================
\section{Extensions}
% ============================================================

SciVault-QF is extensible while maintaining strict core validation.

Application-specific data shall be placed inside an
\texttt{extensions} object.

Example:

\begin{lstlisting}
"extensions": {
  "org.scivault.peer_instruction": {
    "recommended_discussion_time": 120
  }
}
\end{lstlisting}

Reverse-domain namespaces are recommended to reduce naming collisions.

Applications should preserve extension namespaces they do not understand
whenever practical.

An application may ignore unsupported extension data.

Core question correctness should not depend upon proprietary extension data
when portability is intended.

% ============================================================
\section{Separation of Content and Operational Data}
% ============================================================

SciVault-QF represents question content.

It shall not contain identifiable student or session information such as:

\begin{itemize}
    \item student names;
    \item student IDs;
    \item class rosters;
    \item individual responses;
    \item grades;
    \item attendance;
    \item assessment sessions;
    \item peer-instruction rounds; or
    \item response timestamps.
\end{itemize}

Applications shall store operational data separately and reference
SciVault-QF questions using stable identifiers.

% ============================================================
\section{Validation}
% ============================================================

SciVault-QF validation consists of:

\begin{enumerate}
    \item structural validation using the official JSON Schema;
    \item semantic validation using a SciVault-QF validator.
\end{enumerate}

The official schema uses JSON Schema Draft 2020-12.

% ------------------------------------------------------------
\subsection{Validation Severity}
% ------------------------------------------------------------

Findings have one of three severities.

\begin{longtable}{|>{\ttfamily}p{2.5cm}|p{11cm}|}
\hline
error &
The content violates a requirement or cannot be interpreted reliably.
One or more errors make the document invalid. \\ \hline

warning &
The content remains valid but contains a potentially incomplete, ambiguous,
undesirable, or unusual condition. \\ \hline

info &
Informational feedback that does not indicate a defect. \\ \hline
\end{longtable}

A document is \texttt{valid} when it contains no errors.

A document is \texttt{invalid} when it contains one or more errors.

Warnings do not make a document invalid.

% ------------------------------------------------------------
\subsection{Structural Errors}
% ------------------------------------------------------------

The official JSON Schema validates:

\begin{itemize}
    \item required fields;
    \item field data types;
    \item question types;
    \item answer structures;
    \item difficulty values;
    \item media structures;
    \item dynamic-variable structures;
    \item ISO~8601 timestamps; and
    \item unknown core fields.
\end{itemize}

JSON Schema validation failures are errors.

% ------------------------------------------------------------
\subsection{Semantic Errors}
% ------------------------------------------------------------

The semantic validator shall detect conditions including:

\begin{itemize}
    \item duplicate question IDs;
    \item duplicate choice IDs;
    \item answers referencing nonexistent choices;
    \item malformed true/false choices;
    \item variable minimum greater than maximum;
    \item undefined variable placeholders;
    \item undefined variables in expressions;
    \item forbidden expression syntax;
    \item missing required local media;
    \item fill-in-the-blank questions without \texttt{\{\{blank\}\}}; and
    \item violations of mandatory identifier rules.
\end{itemize}

% ------------------------------------------------------------
\subsection{Warnings}
% ------------------------------------------------------------

Warnings may include:

\begin{itemize}
    \item missing explanation;
    \item missing difficulty;
    \item missing tags;
    \item missing objectives;
    \item inadequate image alternative text;
    \item duplicate displayed choice text;
    \item suspiciously similar answer choices;
    \item unusually long question content;
    \item unused dynamic variables;
    \item unused media;
    \item missing question-set ID;
    \item missing question-set version; and
    \item nonstandard extension namespaces.
\end{itemize}

% ------------------------------------------------------------
\subsection{Validation Findings}
% ------------------------------------------------------------

Validators should return structured findings.

Example:

\begin{lstlisting}
{
  "severity": "error",
  "code": "ANSWER_UNKNOWN_CHOICE",
  "message":
    "Answer 'E' does not reference an existing choice.",
  "question_id": "q017",
  "path": "questions[16].answer"
}
\end{lstlisting}

Recommended fields are:

\begin{itemize}
    \item \texttt{severity};
    \item \texttt{code};
    \item \texttt{message};
    \item \texttt{question\_id}, when applicable;
    \item \texttt{path}, when applicable.
\end{itemize}

Stable machine-readable validation codes should be used.

Initial codes include:

\begin{lstlisting}
DUPLICATE_QUESTION_ID
DUPLICATE_CHOICE_ID
ANSWER_UNKNOWN_CHOICE
UNDEFINED_VARIABLE
INVALID_VARIABLE_RANGE
INVALID_EXPRESSION
MISSING_MEDIA
MISSING_BLANK
MISSING_EXPLANATION
MISSING_OBJECTIVES
DUPLICATE_CHOICE_TEXT
UNUSED_VARIABLE
UNUSED_MEDIA
\end{lstlisting}

% ============================================================
\section{Unsupported Features}
% ============================================================

A SciVault application is not required to implement every question type or
feature defined by SciVault-QF.

When encountering valid content it does not support, an application should:

\begin{enumerate}
    \item preserve the content;
    \item identify it as unsupported;
    \item avoid corrupting or deleting it; and
    \item continue processing supported content whenever possible.
\end{enumerate}

For example, an early version of SciVault Peer Instruction may support only
\texttt{multiple\_choice} even though the question bank contains other
valid SciVault-QF question types.

% ============================================================
\section{Accessibility}
% ============================================================

Question authors and applications should consider accessibility throughout
the content lifecycle.

At minimum:

\begin{itemize}
    \item meaningful images shall contain alternative text;
    \item correctness shall not depend solely on color;
    \item mathematical expressions should be represented semantically when
    practical;
    \item questions should remain usable with assistive technology when
    practical; and
    \item descriptions should communicate media information required to
    answer a question.
\end{itemize}

% ============================================================
\section{Licensing}
% ============================================================

Each question collection shall specify a content license.

Example:

\begin{lstlisting}
"license": "CC-BY-4.0"
\end{lstlisting}

The license governing SciVault software is independent of the license
governing SciVault-QF content.

Individual media objects may specify different licenses where necessary.

% ============================================================
\section{Question Types Defined by Version 0.1}
% ============================================================

\begin{longtable}{|>{\ttfamily}p{4cm}|p{9.5cm}|}
\hline

multiple\_choice &
Exactly one correct answer selected from discrete choices. \\ \hline

multiple\_select &
One or more correct answers selected from discrete choices. \\ \hline

true\_false &
Binary true/false question. \\ \hline

short\_answer &
Short textual response. \\ \hline

fill\_in\_the\_blank &
One textual blank contained within a prompt. \\ \hline

numerical &
Numerical response with optional tolerance and physical unit. \\ \hline

\end{longtable}

% ============================================================
\section{Future Considerations}
% ============================================================

Potential future features include:

\begin{itemize}
    \item multipart questions;
    \item multiple independently scored blanks;
    \item ranking;
    \item matching;
    \item ordering;
    \item graphical responses;
    \item image-region selection;
    \item vector drawing;
    \item symbolic equation entry;
    \item extended-response questions;
    \item passages and question groups;
    \item random question pools;
    \item conditional questions;
    \item significant-figure checking;
    \item physical-unit equivalence;
    \item richer dynamic-expression syntax;
    \item embedded SVG diagrams;
    \item distractor misconception metadata;
    \item empirical item-difficulty statistics;
    \item discrimination indices;
    \item standards-framework references;
    \item localization and translated variants;
    \item packaged question-bank formats; and
    \item an expanded standard validation-code registry.
\end{itemize}

% ============================================================
\section{Complete Example}
% ============================================================

\begin{lstlisting}
{
  "format": "SciVault-QF",
  "format_version": "0.1",

  "metadata": {
    "id": "org.scivault.physics.kinematics",
    "title": "Kinematics Concept Questions",
    "description":
      "Introductory kinematics questions.",
    "author": "Jane Doe",
    "license": "CC-BY-4.0",
    "version": "1.0",
    "language": "en-US",
    "tags": [
      "physics",
      "kinematics"
    ]
  },

  "questions": [
    {
      "id": "q001",
      "type": "multiple_choice",

      "question":
        "An object starts from rest and accelerates at "
        "$2\\,\\mathrm{m/s^2}$ for $5\\,\\mathrm{s}$. "
        "What is its final velocity?",

      "choices": [
        {
          "id": "A",
          "text": "$5\\,\\mathrm{m/s}$"
        },
        {
          "id": "B",
          "text": "$10\\,\\mathrm{m/s}$"
        },
        {
          "id": "C",
          "text": "$15\\,\\mathrm{m/s}$"
        },
        {
          "id": "D",
          "text": "$20\\,\\mathrm{m/s}$"
        }
      ],

      "answer": "B",

      "explanation":
        "Using $v=v_0+at$, "
        "$v=0+(2)(5)=10\\,\\mathrm{m/s}$.",

      "difficulty": 2,

      "tags": [
        "physics",
        "kinematics",
        "constant-acceleration"
      ],

      "objectives": [
        "KIN-VELOCITY"
      ],

      "settings": {
        "shuffle_choices": true
      }
    }
  ]
}
\end{lstlisting}

% ============================================================
\section{Conformance}
% ============================================================

A document conforms to SciVault-QF 0.1 when:

\begin{enumerate}
    \item it is valid JSON encoded as UTF-8;
    \item it identifies itself as \texttt{SciVault-QF};
    \item its \texttt{format\_version} is \texttt{0.1};
    \item it passes the official SciVault-QF 0.1 JSON Schema; and
    \item it contains no errors identified by the normative SciVault-QF
    semantic validation rules.
\end{enumerate}

Warnings and informational findings do not prevent conformance.

% ============================================================
\section{Summary}
% ============================================================

SciVault-QF provides the common question-content representation for the
SciVault software ecosystem.

Version 0.1 uses UTF-8 JSON and the standard \texttt{.json} extension.

Question and answer-choice IDs are stable identifiers.

Question content is kept separate from student and session data.

SciVault-QF supports scientific and mathematical notation, media,
difficulty metadata, learning objectives, dynamic variables, several
question types, validation, and namespaced extensions.

Future specifications may extend the question model and define a portable
package format for question banks and associated resources.

\end{document}
```